% SPDX-License-Identifier: CC-BY-SA-4.0
% Author: Matthieu Perrin

\begingroup

\begin{exercice}[Automate décidant un langage]\label{exo:lexing/automata/write_automata}

  Proposez un automate fini déterministe reconnaissant chacun des langages ci-dessous.
  \begin{question}
  \item Le langage des mots sur l'alphabet $\{a,b\}$ contenant un nombre de $a$ divisible par $3$. 
  \item Le langage des écritures binaires de nombres entiers divisibles par $4$. 
  \item Le langage des mots sur l'alphabet $\{a,b\}$ contenant au moins un facteur $ab$ et un facteur $ba$. 
  \end{question}

\end{exercice}

\begin{correction}

  \begin{question}
  \item On encode le nombre de $a$ modulo 3 dans les états de l'automate :
    \begin{tikzpicture}[automaton, baseline=(0.base), y=10mm]
      \state[initial, accepting]   (0) at (1,0) {$0$}; 
      \state                       (1) at (0,1) {$1$}; 
      \state                       (2) at (2,1) {$2$}; 

      \path (0) edge[loop right] node {$b$} (0);
      \path (1) edge[loop left]  node {$b$} (1);
      \path (2) edge[loop right] node {$b$} (2);
      \path (0) edge             node {$a$} (1);
      \path (1) edge             node {$a$} (2);
      \path (2) edge             node {$a$} (0);
    \end{tikzpicture}

  \item En se basant sur l'identité $(2n + m) \mod 3 = (2(n \mod 3) + m) \mod 3$,
    on peut n'encoder que le reste de la division par 3 du préfixe déjà analysé, en 3 états (0, 1 et 2).
    Les transitions entre ses trois états sont obtenues grâce à la table de division suivante à gauche.
    Il faut ajouter deux états $I$ et $F$ pour empêcher les $0$ inutiles à gauche du mot et le nombre $0$.

    \begin{minipage}{.5\textwidth}
      \centering
      $\begin{array}{|c||c|c|}
        \hline
        n\mod 3 & 2n+0\mod 3 & 2n+1\mod 3\\
        \hline
        0 & 0 & 1 \\
        1 & 2 & 0 \\
        2 & 1 & 2 \\
        \hline
      \end{array}$
    \end{minipage}%
    \begin{minipage}{.3\textwidth}
      \begin{tikzpicture}[automaton, baseline=(1.base)]
        \state[initial]   (I) at (0,1) {$I$}; 
        \state[accepting] (F) at (0,0) {$F$}; 
        \state[accepting] (0) at (1,0) {$0$}; 
        \state            (1) at (1,1) {$1$}; 
        \state            (2) at (2,1) {$2$}; 

        \path (I) edge             node       {$0$}    (F);
        \path (I) edge             node       {$1$}    (1);
        \path (1) edge[bend left]  node       {$1$}    (0);
        \path (0) edge[bend left]  node       {$1$}    (1);
        \path (1) edge[bend left]  node       {$0$}    (2);
        \path (2) edge[bend left]  node       {$0$}    (1);
        \path (0) edge[loop right]  node       {$0$}    (0);
        \path (2) edge[loop right]  node       {$1$}    (2);
      \end{tikzpicture}
    \end{minipage}%

  \item Le langage est reconnu par l'automate suivant : 
    \begin{tikzpicture}[automaton, baseline=(1.base), y=10mm]
      \state[initial]   (1) at (0,1) {$1$}; 
      \state            (2) at (1,2) {$2$}; 
      \state            (3) at (1,0) {$3$}; 
      \state            (4) at (2,2) {$4$}; 
      \state            (5) at (2,0) {$5$}; 
      \state[accepting] (6) at (3,1) {$6$}; 

      \path (2) edge[loop left]  node       {$a$}    (2);
      \path (3) edge[loop left]  node       {$b$}    (3);
      \path (4) edge[loop right] node       {$b$}    (4);
      \path (5) edge[loop right] node       {$a$}    (5);
      \path (6) edge[loop right] node       {$a, b$} (6);
      \path (1) edge             node[swap] {$a$}    (2);
      \path (1) edge             node       {$b$}    (3);
      \path (2) edge             node[swap] {$b$}    (4);
      \path (3) edge             node       {$a$}    (5);
      \path (4) edge             node[swap] {$a$}    (6);
      \path (5) edge             node       {$b$}    (6);
    \end{tikzpicture}
    
  \end{question}

\end{correction}

\endgroup
\endinput
